% Loi de décroissance radioactive : nombre de noyaux restants après des demi-vies successives (L'essentiel, chapitre 6)
\begin{tikzpicture}
\begin{axis}[otMesures,otAxesBout,width=104mm,height=40mm,xmin=0,xmax=5.3,ymin=0,ymax=1.1,
  xtick={0,1,...,5},xticklabels={$0$,$t_{1/2}$,$2\,t_{1/2}$,$3\,t_{1/2}$,$4\,t_{1/2}$,$5\,t_{1/2}$},
  ytick={0,0.125,0.25,0.5,1},yticklabels={$0$,$N_0/8$,$N_0/4$,$N_0/2$,$N_0$},
  xlabel={$t$},ylabel={$N(t)$}]
\addplot[otNavy,line width=1pt,domain=0:5.2,samples=120] {2^(-x)};
\foreach \k/\v in {1/0.5,2/0.25,3/0.125}{%
  \edef\temp{\noexpand\draw[otOrange,line width=0.6pt,dash pattern=on 2pt off 1.5pt] (axis cs:\k,\v) -- (axis cs:\k,0) (axis cs:\k,\v) -- (axis cs:0,\v);}\temp
  \edef\temp{\noexpand\addplot[only marks,mark=*,mark size=1.8pt,otOrange] coordinates {(\k,\v)};}\temp}
\node[pcLegende,anchor=north east,text=otNavy] at (axis cs:5.2,1.05) {$N(t)=N_0\,\mathrm{e}^{-\lambda t}$, \quad $t_{1/2}=\dfrac{\ln2}{\lambda}$};
\end{axis}
\end{tikzpicture}
\par\smallskip\legende{Au bout de $n$ demi-vies, il reste $N_0/2^n$ noyaux : la moitié des noyaux restants se désintègre pendant chaque durée $t_{1/2}$.}
